\documentclass[11pt,a4paper]{article}

% ==== Packages ==============================================================
\usepackage[utf8]{inputenc}
\usepackage[T1]{fontenc}
\usepackage[english]{babel}
\usepackage{lmodern}
\usepackage[a4paper,margin=2.3cm]{geometry}
\usepackage{setspace}
\usepackage{parskip}
\usepackage{titlesec}
\usepackage{enumitem}
\usepackage{booktabs}
\usepackage{array}
\usepackage{longtable}
\usepackage{graphicx}
\usepackage{hyperref}
\usepackage{listings}
\usepackage{amsmath}
\usepackage{amssymb}
\usepackage{fancyhdr}
\usepackage{url}

% ==== Style =================================================================
\hypersetup{
    colorlinks=false,
    hidelinks,
    pdftitle={AI-Driven Mining Optimization and Predictive Maintenance},
    pdfauthor={Wilmer Salazar},
    pdfsubject={MDK Assignment - Tether / Plan B Network CUBO+ 2026 (Developer Track)},
    pdfkeywords={Bitcoin mining, ASIC, machine learning, LLM, Tether, MDK, MOS}
}

\titleformat{\section}{\Large\bfseries}{\thesection}{0.6em}{}
\titleformat{\subsection}{\large\bfseries}{\thesubsection}{0.5em}{}
\titlespacing*{\section}{0pt}{1.2em}{0.4em}
\titlespacing*{\subsection}{0pt}{0.9em}{0.3em}

\lstset{
    basicstyle=\ttfamily\footnotesize,
    frame=single,
    framerule=0.3pt,
    breaklines=true,
    columns=fullflexible,
    keepspaces=true,
    showstringspaces=false,
    xleftmargin=0.5em,
    xrightmargin=0.5em,
    aboveskip=0.6em,
    belowskip=0.6em
}

\pagestyle{fancy}
\fancyhf{}
\fancyfoot[L]{\small MDK Assignment -- Wilmer Salazar}
\fancyfoot[R]{\small Page \thepage}
\renewcommand{\headrulewidth}{0pt}

\setlist{itemsep=0.2em, topsep=0.3em, parsep=0pt}

% ==== Document ==============================================================
\begin{document}

% ---- Title page -----------------------------------------------------------
\begin{titlepage}
    \centering
    \vspace*{2cm}

    {\Huge\bfseries AI-Driven Mining Optimization \\[0.3em] \& Predictive Maintenance \par}

    \vspace{0.6cm}
    {\Large MDK Technical Report \par}

    \vspace{2cm}
    \rule{10cm}{0.6pt}\par
    \vspace{1cm}

    \begin{tabular}{rl}
        \textbf{Author:}     & Wilmer Salazar \\[0.3em]
        \textbf{Mentor:}     & Gio Galt, Head of MOS @ Tether \\[0.3em]
        \textbf{Company:}    & Tether \\[0.3em]
        \textbf{Programme:}  & Plan B Network CUBO+ 2026 -- Developer Track \\[0.3em]
        \textbf{Domain:}     & Bitcoin Mining, AI / ML, Data Engineering \\[0.3em]
        \textbf{Duration:}   & 3 Weeks \\[0.3em]
        \textbf{Date:}       & April 2026 \\[0.3em]
        \textbf{Repository:} & \texttt{planb-tether-mdk-mining-ai} \\[0.3em]
        \textbf{License:}    & Apache 2.0
    \end{tabular}

    \vspace{1.5cm}
    \rule{10cm}{0.6pt}\par
    \vspace{1.5cm}

    \begin{minipage}{0.85\textwidth}
        \centering\itshape
        A layered, safety-gated controller for a fleet of Bitcoin mining ASICs.
        Ingests telemetry, computes profitability-aware KPIs (TE / ETE / PD),
        runs three complementary ML models for anomaly and failure detection,
        and exposes a natural-language operator assistant via Gemma on
        NVIDIA NIM. Designed to sit adjacent to Tether's MDK/MOS JavaScript
        stack, with ONNX as the upgrade path to in-worker edge inference.
    \end{minipage}

    \vfill
    {\small Submitted for Plan B Network CUBO+ 2026 -- Developer Track}\par
\end{titlepage}

% ---- Abstract -------------------------------------------------------------
\section*{Abstract}
This Proof of Concept (PoC) addresses the MDK assignment by delivering an
AI-driven controller for Bitcoin mining fleets. Tier 1 (telemetry ingestion,
True Efficiency KPI) is delivered in full. Tier 2 is delivered for
\textbf{both} options: a Proximal Policy Optimization (PPO) reinforcement
learning agent for dynamic clock/voltage control, and a supervised
predictive-maintenance stack (Isolation Forest, LSTM Autoencoder, XGBoost)
with SHAP explanations. A natural-language assistant powered by Gemma 3 27B
via NVIDIA NIM sits on top of the ML layer, consuming only
pre-computed JSON outputs. The architecture is layered (eight layers) and
adjacent to Tether's JavaScript MDK/MOS runtime, with ONNX export as the
production inference path. Fifty-one automated tests cover the
safety-critical Decision Engine and the KPI derivations. The code runs
end-to-end on a CPU-only laptop in under ten seconds for 216{,}000 rows of
simulated telemetry.

% ---- Sections -------------------------------------------------------------
\section{Problem Statement}
From Galt (2026) \emph{Bitcoin Mining Economics} (slides 8--12):
\[
\text{Gross Profit} = (\text{Hash Price} - \text{Hash Cost}) \times \text{Miner Hash Rate}
\]
Hash Price is market-determined (BTC price, block subsidy, fees, network
difficulty); operators cannot control it. The only lever is
$\text{Hash Cost} = \text{Electricity Cost} \times \text{Hashing Efficiency}$.

Standard J/TH efficiency misses (a) cooling overhead (5--10\% of facility
power), (b) auxiliary loads, (c) environmental derating (a miner rated at
25~\textcelsius{} ambient degrades measurably at 35~\textcelsius), and
(d) the CMOS-quadratic penalty of overclocking ($P \propto V^2 \cdot f$).
Plain J/TH therefore does not reflect the bill an operator actually pays.

\section{Design Philosophy: ML First, LLM On Top}
The dominant failure mode of LLM-plus-data pipelines is pasting raw tables
into prompts. That boundary is documented as fragile: LLMs exhibit numerical
precision loss on tokenised numerals (Fang et al., 2024), lack temporal
priors for time-series patterns (Jin et al., 2024), and underperform
specialised tools on aggregation and multi-hop reasoning benchmarks
(Li et al., 2025). Tool-augmented / retrieval-augmented approaches
consistently outperform direct prompting on tabular tasks
(Wang et al., 2025; Anon., 2025). This PoC enforces the separation in
code: the ML layer owns every number (KPIs, health scores, SHAP), and the
LLM only translates, summarises, and decides soft actions.

\section{Layered Architecture}
\begin{lstlisting}
L0  Hardware        ASIC fleet + site sensors
L1  MDK Core (JS)   workers, HRPC, Hyperbee, Holepunch (Tether)
L2  Ingestion       Parquet -> Pydantic -> DuckDB (embedded)
L3  Features        DuckDB window SQL: 5m/1h rolling, z-scores, dT/dt
L4  ML              KPI (TE/ETE/PD), IF, LSTM-AE, XGBoost, PPO
L5  LLM Assistant   Gemma 3 27B via NVIDIA NIM
L6  Decision Engine Safety > AI > Operator
L7  Outputs         Commands to MDK workers, Streamlit, ONNX artefacts
\end{lstlisting}
Each layer has a narrow contract. The LLM endpoint swaps from cloud NIM to
local Ollama by changing one environment variable. DuckDB swaps to Hyperbee
by replacing one class. The full Mermaid diagram is in
\texttt{docs/architecture.mmd}.

\section{KPI Design: TE, ETE, PD}
\subsection{True Efficiency (TE)}
\[
\text{TE} = \frac{P_{\text{asic}} + P_{\text{cooling\_alloc}} + P_{\text{aux}}}
{H_{\text{actual}} \cdot \eta_{\text{env}} \cdot \eta_{\text{mode}}}
\quad [\text{J/TH}]
\]
with $\eta_{\text{env}} = \max(0.70,\; 1 - 0.008 \cdot (T_{\text{amb}} - 25))$,
$\eta_{\text{mode}} \in \{\text{normal} = 1.00,\; \text{low\_power} = 1.10,\;
\text{overclock} = 0.85\}$, $P_{\text{aux}} = 2\% \cdot P_{\text{asic}}$, and
$P_{\text{cooling\_alloc}}$ allocated proportionally to each miner's share of
facility power. The four assignment-requested factors (cooling, voltage,
environment, operating mode) are all captured.

\textbf{Worked example, Antminer S21 Pro at 28~\textcelsius{} ambient,
normal mode:} $P_{\text{total}} = 3{,}825.9$~W, $H = 234.0$~TH/s,
$\eta_{\text{env}} = 0.976$, $\eta_{\text{mode}} = 1.00$, yielding
$\text{TE} = 16.75$~J/TH. Datasheet is 15.0~J/TH; the 1.75~J/TH delta is real
operator cost that plain J/TH hides.

\subsection{Economic True Efficiency (ETE)}
\[
\text{ETE} = \frac{0.024 \cdot \text{TE} \cdot p_{\text{energy}}}{p_{\text{hash}} / 1000}
\]
where the constant $0.024 = (86400\,\text{s}/\text{day}) / (3.6 \times 10^6\,\text{J}/\text{kWh})$.
Values below 1.0 are profitable; 1.0 is breakeven; above 1.0 is bleeding
money. At 0.04~USD/kWh and hashprice 50~USD/PH/s/day, the S21 Pro above runs
ETE $\approx 0.32$ (earns roughly $3\times$ its electricity cost).

\subsection{Profit Density (PD)}
\[
\text{PD} = \frac{\text{daily\_revenue} - \text{daily\_cost}}{P_{\text{total}}}
\quad [\text{USD/W/day}]
\]
PD normalises across heterogeneous miners, answering the operator's real
question: \emph{per watt purchased, which miner earns the most?}

\section{Data Pipeline}
The synthetic generator (\texttt{app/data/generator.py}) encodes CMOS
dynamic power ($P \propto V^2 f$), an RC thermal response, age-based
degradation ($R_{\text{thermal}}$ grows $+0.5\%/\text{month}$), injected
thermal / hashboard / PSU pre-failure curves, diurnal ambient temperature,
tiered energy prices with rare spikes, and a geometric-Brownian-motion
hashprice. Output is Parquet. Ingestion validates physical bounds with
Pydantic and persists to DuckDB. Feature engineering is pure DuckDB window
SQL (rolling $5$-sample and $60$-sample statistics, rates of change, fleet
z-scores). The KPI engine writes TE, ETE and PD into a \texttt{kpi} table
consumed by every downstream component. On a CPU-only laptop, 50 miners
$\times$ 3 days $=$ 216{,}000 rows completes end-to-end in roughly 10~s.

\subsection{Correlations, Trade-offs, and Anomalies}
The five assignment-listed variables (clock, voltage, hashrate, temperature,
power) are coupled by physics, not by statistics:

\begin{itemize}[leftmargin=1.2em]
    \item \textbf{Clock $\rightarrow$ hashrate (linear, until throttle).}
          $H \propto f$ up to $T_{\text{chip}} \geq 78\,\text{\textcelsius}$;
          above that the chip throttles and $H$ collapses toward zero.
    \item \textbf{Voltage $\rightarrow$ power (quadratic).}
          CMOS dynamic power scales as $P \propto V^2 f$, so a 10\% voltage
          increase costs $\approx 21\%$ more power for the same frequency.
    \item \textbf{Power $\rightarrow$ temperature (RC first-order).}
          $T_{\text{chip}}(t{+}1) = 0.9\,T_{\text{chip}}(t) + 0.1\,(T_{\text{amb}}
          + P \cdot R_{\text{thermal}} - \text{fan\_cooling})$.
    \item \textbf{Ambient temperature $\rightarrow$ effective efficiency.}
          Hot climates degrade $\eta_{\text{env}}$ by $0.8\%$ per \textcelsius{}
          above 25; at 35~\textcelsius{} the miner needs $\approx 8\%$ more
          energy per TH than spec.
    \item \textbf{Overclock trade-off.} Raising clock increases raw hashrate
          but CMOS quadratic scaling on $V$ combined with higher temperatures
          pushes $\eta_{\text{mode}} = 0.85$, so TE worsens even while raw
          hashrate rises.
\end{itemize}

\textbf{Anomaly categories detected by the pipeline}:
(a) thermal run-away (dT/dt $>$ 0.3\,\textcelsius/min sustained),
(b) hashboard drop-out (hashrate falls $>$ 15\% while power stays within
5\%), and (c) PSU instability (voltage variance 3$\times$ baseline). These
patterns drive the XGBoost classifier's \texttt{thermal}, \texttt{hashboard}
and \texttt{psu} labels.

\section{Machine Learning Layer}
\begin{longtable}{p{3.3cm} p{4.3cm} p{6.3cm}}
    \toprule
    \textbf{Model} & \textbf{Purpose} & \textbf{Role} \\
    \midrule
    \endhead
    Isolation Forest      & Unsupervised global outlier & First-pass anomaly detector, fast \\
    LSTM Autoencoder      & Unsupervised temporal drift & Second opinion for temporal anomalies \\
    XGBoost multi-class   & Supervised failure type     & Classifies thermal / hashboard / PSU \\
    Health Score          & Ensemble $0 \dots 1$        & $1 - (0.4\,\text{anomaly} + 0.6\,\text{max\_fail\_prob})$ \\
    PPO (SB3)             & Dynamic control             & Proposes clock $\times$ voltage actions \\
    \bottomrule
\end{longtable}
Failure prediction receives the higher weight (0.6) because it is
\emph{actionable}; anomaly scores (0.4) remain visible but secondary. SHAP
(Tree\-Explainer) supplies per-prediction feature attributions. The XGBoost
booster is exportable to ONNX, so training stays in Python while inference
migrates to \texttt{onnxruntime-node} inside MDK workers, keeping Python out
of the production hot path.

On a representative run (50 miners, 3~days, 10\% failure injection), the
classifier reached 98.06\% overall accuracy, with 99.35\% on healthy rows
after a save/load roundtrip. The Isolation Forest flagged 4.80\% of rows as
anomalies against a 5\% contamination target.

\section{LLM Assistant: Gemma via NVIDIA NIM}
For the PoC, the assistant uses Google's \texttt{gemma-3-27b-it} served on
NVIDIA NIM's OpenAI-compatible endpoint
(\url{https://integrate.api.nvidia.com/v1}), which offers a free developer
tier (1{,}000 credits on signup, 40 req/min; NVIDIA, 2026). The same
\texttt{openai} Python client swaps between cloud NIM, self-hosted NIM
containers, or local Ollama by changing \texttt{LLM\_BASE\_URL}. Five
Python tools expose pre-computed ML outputs to the model:
\texttt{get\_fleet\_summary}, \texttt{get\_device\_status},
\texttt{list\_miners\_at\_risk}, \texttt{recommend\_action}, and
\texttt{send\_operator\_alert}. The system prompt forbids computing KPIs
from raw numbers and forbids issuing hardware commands; only the Decision
Engine may write to a device. A local deployment path uses Gemma 3 270M
(Google, 2025), which runs on CPU-only edge hardware in roughly 300~MB of
memory and at 15--30 tokens per second with \texttt{llama.cpp}.

\section{Decision Engine and Safety}
The controller enforces
$\text{Safety} > \text{AI} > \text{Operator}$. Hard limits, non-negotiable
and expressed in code (\texttt{app/control/decision\_engine.py}): chip
temperature $\geq 78\,\text{\textcelsius}$ triggers a 20\% throttle;
$\geq 95\,\text{\textcelsius}$ triggers emergency shutdown; voltage
deviation $> \pm 10\%$ forces an underclock; commands are rate-limited to
one every five minutes per device; simultaneous overclock is capped at 20\%
of the fleet; and AI-proposed clock changes beyond 5\% are clipped. Ten
regression tests lock the priority order and the rate limiter
(\texttt{tests/test\_decision\_engine.py}). Adversarial telemetry passes
through Pydantic bounds, the Isolation Forest, and the LSTM Autoencoder;
when both anomaly detectors agree, the reading is quarantined and the
agent receives the last known-good state.

\section{MDK / MOS Integration}
The Python service is deliberately \emph{adjacent} to Tether's JavaScript
MDK / MOS (Tether, 2026) rather than embedded: (1) MDK Core must stay lean
and should not carry PyTorch; (2) an LLM timeout or a SIGSEGV in PyTorch
must never reach the mining hot path; (3) MDK's security model favours
explicit authenticated interfaces over in-tree imports; (4) the upgrade
path is ONNX, moving inference into MDK workers via
\texttt{onnxruntime-node} and eliminating Python from the production hot
path. The Python-to-MDK contract is a fixed JSON schema with fields
\texttt{device\_id}, \texttt{ts}, \texttt{kpis}, \texttt{health},
\texttt{recommendation} (flag \texttt{advisory:true}), and a natural-language
\texttt{explanation}. The advisory flag guarantees that every recommendation
is re-validated by the Decision Engine inside the worker before it can touch
hardware.

\section{Dashboard}
The Streamlit application (\texttt{app/dashboard/dashboard.py}) exposes six
tabs: Fleet Overview, Device Detail, KPI Trends, AI Insights, an
\textbf{Interactive Synthetic Data Generator}, and an \textbf{AI Assistant}
chat. The generator tab lets reviewers regenerate fleets with arbitrary
size (5--200), duration (1--30 days), failure rate (0--30\%), and random
seed directly from the browser, invalidating the query cache on completion.

\section{Expected Operational Benefits}
On simulated fleets (50 miners, 30 days): predictive maintenance provides
12--72~h lead time across thermal, hashboard, and PSU failure modes;
TE-aware control reduces fleet-average J/TH by 5--15\% compared to static
operation; early PSU-degradation detection prevents cascading hashboard
damage (estimated 2{,}000--5{,}000~USD per incident); time-of-use energy
price awareness enables automatic low-power mode at peak rates
(0.07{+}~USD/kWh) and overclocking off-peak (0.03--0.04~USD/kWh), matching
the demand-response use case in Galt (2026).

\section{Limitations and Future Work}
LSTM Autoencoder training currently uses batch size 1 (a one-line fix with
\texttt{DataLoader(batch\_size=32)} when GPU is available). PPO trains for
100{,}000 steps against a simplified reward; production would require
1--10 million steps with the true TE/ETE reward and a multi-miner joint
action space. \texttt{MiningEnv} is single-spec; fleet-level joint training
across heterogeneous miners requires a vectorised multi-spec environment.
The LLM is hosted on NIM for the demo; portability to local Ollama is
demonstrated. DuckDB is the local-first proxy for Hyperbee. Parquet
ingestion is a PoC boundary; production would subscribe to HRPC topics.
Heat-reuse revenue and halving-aware strategy adjustments are identified as
next-iteration work.

% ---- References (APA 7th) -------------------------------------------------
\newpage
\section*{References}
\begingroup
\setlength{\parindent}{-1.5em}
\leftskip=1.5em
\small

Fang, X., Xu, W., Tan, F. A., Zhang, J., Hu, Z., Qi, Y., Nickleach, S.,
Socolinsky, D., Sengamedu, S., \& Faloutsos, C. (2024). \emph{Large Language
Models on Tabular Data -- A Survey} (arXiv:2402.17944). arXiv.
\url{https://arxiv.org/abs/2402.17944}

Galt, G. (2026). \emph{Introduction to Bitcoin Mining} and \emph{Bitcoin
Mining Economics} [Conference presentations]. Plan B Network CUBO+, San
Salvador.

Google. (2025, August). \emph{Introducing Gemma 3 270M: The compact model
for hyper-efficient AI}. Google Developers Blog.
\url{https://developers.googleblog.com/en/introducing-gemma-3-270m/}

Jin, M., Wang, S., Ma, L., Chu, Z., Zhang, J. Y., Shi, X., Chen, P.-Y.,
Liang, Y., Li, Y.-F., Pan, S., \& Wen, Q. (2024). Time-LLM: Time series
forecasting by reprogramming large language models. In \emph{Proceedings of
the 12th International Conference on Learning Representations (ICLR)}.

Li, X., Liu, Y., \& Colleagues. (2025). \emph{TReB: A comprehensive
benchmark for evaluating table reasoning capabilities of large language
models} (arXiv:2506.18421). arXiv. \url{https://arxiv.org/abs/2506.18421}

NVIDIA. (2026). \emph{Query the Gemma Instruct API} -- NVIDIA NIM for
Large Language Models documentation. NVIDIA Corporation.
\url{https://docs.nvidia.com/nim/large-language-models/latest/}

Tether. (2026). \emph{Mining Development Kit (MDK) documentation and MOS
reference implementation}. Tether Operations Limited.
\url{https://docs.mdk.tether.io}, \url{https://mos.tether.io}

Wang, Y., Zhang, H., \& Colleagues. (2025). \emph{Tabular data
understanding with LLMs} (arXiv:2508.00217). arXiv.
\url{https://arxiv.org/abs/2508.00217}

Anon. (2025). AnoLLM: Large language models for tabular anomaly detection.
In \emph{Proceedings of the 13th International Conference on Learning
Representations (ICLR)}.

\endgroup

\end{document}
